\documentclass[11pt]{article}
\usepackage[margin=1in]{geometry}
\usepackage{amsmath,amssymb,amsthm}
\usepackage{graphicx}
\usepackage[hidelinks]{hyperref}

\newtheorem{theorem}{Theorem}
\newtheorem{lemma}{Lemma}

\newcommand{\Cp}{C_p}
\newcommand{\supp}{\operatorname{supp}}
\newcommand{\spanlin}{\operatorname{span}}

\title{A Convergent-Sequence Subcase Of A Weak-Star Basic Sequence Problem}
\author{FA/Banach smoke run packet}
\date{June 9, 2026}

\begin{document}
\maketitle

\section*{Source And Target}

The source paper is Jerzy Kakol, Manuel Lopez-Pellicer, and Wieslaw Sliwa,
\emph{On weak*-basic sequences in duals and biduals of spaces C(X) and
Quojections}, arXiv:2601.19873.

The target is Problem 12 of the paper:
\begin{quote}
Does $\Cp(X)^*$ admit a weak$^*$-basic sequence for any infinite compact
space $X$?
\end{quote}

The screenshot below is a full-width crop from page 7 of the source PDF.

\begin{center}
\includegraphics[width=\textwidth]{figures/open_problem_crop.png}
\end{center}

\section*{Candidate Partial Result}

\begin{theorem}
Let $X$ be a compact Hausdorff space containing a nontrivial convergent
sequence. Then $\Cp(X)^*$ admits a weak$^*$-basic sequence.
\end{theorem}

\begin{proof}
Let $(x_n)$ be a sequence of pairwise distinct points of $X$ converging to a
point $x\in X\setminus\{x_n:n\in\mathbb N\}$. For each $n$, define
\[
  \mu_n=\delta_{x_n}-\delta_x\in \Cp(X)^*.
\]
Here $\delta_t$ denotes point evaluation at $t$. Since the topology of
$\Cp(X)$ is pointwise convergence, every finite linear combination of point
evaluations is continuous.

Let $K=\{x\}\cup\{x_n:n\in\mathbb N\}$. This set is compact and closed in
$X$. We first show that
\[
  M:=\spanlin\{\mu_n:n\in\mathbb N\}
\]
is weak$^*$ closed in $\Cp(X)^*$, where the weak$^*$ topology is
$\sigma(\Cp(X)^*,\Cp(X))$.

Take $\nu\in\overline{M}^{\,w^*}$. Since the continuous dual of $\Cp(X)$
consists of finitely supported signed measures, write
\[
  \nu=\sum_{j=1}^m a_j\delta_{t_j}.
\]
If some $t_j\notin K$ has nonzero coefficient, normality of the compact
Hausdorff space $X$ gives a function $f\in C(X)$ such that $f(t_j)=1$ and
$f$ vanishes on $K$ and on the remaining finitely many support points of
$\nu$. Every element of $M$ vanishes on this $f$, so $\nu(f)=0$, a
contradiction. Thus $\supp\nu\subset K$.

The support of $\nu$ is finite, hence
\[
  \nu=a\delta_x+\sum_{n=1}^N b_n\delta_{x_n}
\]
for some $N$. Every $\mu\in M$ annihilates the constant function $1$, so
the same is true of $\nu$ by weak$^*$ closure. Therefore
$a+\sum_{n=1}^N b_n=0$, and consequently
\[
  \nu=\sum_{n=1}^N b_n(\delta_{x_n}-\delta_x)
      =\sum_{n=1}^N b_n\mu_n\in M.
\]
So $M$ is weak$^*$ closed.

It remains to check that $(\mu_n)$ is a Schauder basis for $M$ with the
subspace weak$^*$ topology. Each element of $M$ has a unique finite expansion
in the sequence $(\mu_n)$. For each fixed $n$, the point $x_n$ is isolated in
the compact set $K$, so $K\setminus\{x_n\}$ is closed in $X$. By Urysohn's
lemma there is $f_n\in C(X)$ such that
\[
  f_n(x_n)=1,\qquad f_n\vert_{K\setminus\{x_n\}}=0.
\]
Then, for every $\sum_k c_k\mu_k\in M$,
\[
  \left(\sum_k c_k\mu_k\right)(f_n)=c_n.
\]
Thus the $n$th coefficient functional is weak$^*$ continuous on $M$. Since
all expansions in $M$ are finite, the partial sums converge trivially in
the subspace topology. Hence $(\mu_n)$ is a weak$^*$-basic sequence in
$\Cp(X)^*$.
\end{proof}

\section*{Verification Notes}

The proof uses only standard facts: the continuous dual of $\Cp(X)$ consists
of finitely supported point evaluations, compact Hausdorff spaces are normal,
and a nontrivial convergent sequence plus its limit is closed compact.

The key verification point is the closed-span argument. The total-mass
condition comes from evaluating at the constant function $1$, while support
outside the convergent compactum is excluded by Urysohn separation.

\section*{Limitations And Novelty Check}

This is a partial result only. It does not address compact spaces with no
nontrivial convergent sequence.

There is also a substantial novelty caveat. The source paper later proves a
metrizable-quotient theorem for $\Cp(X)$; after observing that restriction to
the convergent compact subset is an open quotient onto an infinite-dimensional
metrizable locally convex space, this subcase appears to be implicit in the
paper's own methods. The value of this packet is therefore the direct
constructive proof, not a claim of a full new answer to Problem 12.

\section*{Human Review Recommendation}

Likely valid as a constructive partial result. Do not count it as a full solve
of Problem 12.

\end{document}
